\documentclass[11pt]{article}
\usepackage{amsmath,hyperref}
\title{A carried order's settlement includes delivery of its actual completion}
\begin{document}\maketitle
\paragraph{Lift, obligation and envelope.}
When a crew inherits an old warehouse order, a loss of the next
carrier leaves both a station duty and a deadline for the result
to reach the controller. Protect the exact ordered completion,
its named report at the consuming callback, and a separate
local unload. A productive optional job takes one shared pickup
and sends its own 60-byte result to that same consumer. It is
not mandated outside the extra admitted envelope.
The native fragments are the already pinned warehouse station,
inventory and scheduler callbacks and the SubT direct
broker/byte-window/application callbacks. The combined task
uses a stipulated work-to-report adapter, not an existing real
warehouse radio integration. In particular, the scope assumes
fixed authorized senders, the current episode binding, link
availability at relevant cuts, and no RF erasure or unmodeled
competitor. It does not establish those prerequisites for a
real network or an independent mission.

\paragraph{Same history, not a product of pass flags.}
After the initial A loss, classes 7 and 6 are repicked. The
still-active salvage may or may not have moved class 4 to a
named buffer before B is lost. At the actual native buffer read
at tick 9 the unchanged admission core's action table is
queried using the observed contents or the erased read, not
the phase's latent branch. The optional stage if approved
actually puts class 8 in an independent live carrier and
charges one of the shared pick opportunities. The old
station can complete at 12 from its named remaining class
sources. The wrapper sends the 173-byte old report at tick 13
only if the current native station so records, with its
emitted result as the body. The SubT application callback must
return that body; native TX alone counts as no receipt.
No callback is generated for the old report if the old order
never completed in this episode. Separate Home unloading
actually runs even on the negative controls. The radio clock
in the fixture is one quarter second per work tick.
The model has no radio-to-pick feedback over this window,
so the replay can compute the warehouse log and then pass
its actions to the radio in their modeled chronological order.

\paragraph{Compositional budget law.}
There are two separate authorities: use one pickup for the new
job, then consume the named receiver's time-window capacity.
A later closure of the first resource is not an earlier
release of the second. Let \(P\) be total pickup cap and \(C\)
be that receiver's capacity in the window. The old order
requires \(p=3\) pickups with a positive qualified buffer
read and \(p=4\) otherwise. Without a read, the available
and unavailable buffer branches are both in the same
receiving class. If the optional action is taken, it must
fulfill both actual operations. Its capacity consumption is
one pickup and 60 received bytes at tick 11; the old 173
at tick 13 is within the same inclusive window. The exact
positive joint condition for staging early is
\[
                  P\ge p+1,\qquad C\ge233 .
\]
The conditions are necessary in the declared fixture: in
the missing-custody branch the unavailable item must actually
be repicked; a successful earlier consume remains live at
13. They suffice by the original ordered buffer/repick
execution and by the direct admitted broker. Baseline old
operation without the optional job requires \(P\ge p\)
and \(C\ge173\). There are no new cross-product hypotheses
about the joint behavior of two unrelated probes: the native
warehouse completion is literally the only old datum
forwarded by the native callback chain.

A pick-only policy can stage the expected new work and
finish the old order but strand the required old report.
An RX-only policy can fit the old and new packet sizes
but spend the pickup that leaves no actual order to report.
An old failure after a sending time still needs its input
read and source qualification; erasing a failed TX/RX
episode from the log is not a lawful repair. On the matched
144-run sweep, the combined check preserves both old
mandatory outcomes and the native unload in all 32 positive
\(P\ge4\) cases, and permits early new work in five.
The pick-only control loses 15 old receipts after actual
station completion. The RX-only control foregoes seven old
station completions when it stages the new job. Most of
these are sufficiency controls for missing premises, not
matched performance competitors. Replacing the joint test
by blanket refusal of the extra job is safe in the positive
envelope but sacrifices those reachable extra outcomes.

\paragraph{Baselines and remaining gap.}
A typed two-budget scheduler with the same source receipt,
actual read, deadlines and capacity ties the returned
decisions; callback effects were checked separately from
the arithmetic. \href{https://lillyclark.github.io/files/ACHORD_RA_L_IROS_.pdf}{ACHORD} already
exposes queue data to planning, distinguishes message classes
and uses acknowledgement and configurable rate allocation.
Those are serious components for an alternative assembly;
ordinary after-enqueue priority does not by itself reserve
space for a \emph{future} old body while it does not yet
exist. With a fixed reserve for this one body and the same
effect link, a direct assembly absorbs the example.
\href{https://arxiv.org/abs/2103.11470}{NeBula}'s relay and
return strategies address a path loss that is held fixed
here, not compared as executed traces. This is positive
scoped evidence of active physical work and transport
on a common modeled history, not a result against the full
established architecture, support for every fault inventory
row or a held-out transfer. Variable paths, two live receipt
owners or hidden side effects on the read require new
source-qualified premises.
\end{document}
